%==================================================================
\section{Complément de la nuit du 22 au 23 septembre 2026}\label{sec:nuit}
\emph{Travail autonome (l'opérateur dormait). Même échelle de statut. Détail horodaté : \code{NUIT\_JOURNAL.md}.}

\subsection{Recherche bibliographique ciblée}
\begin{itemize}[nosep]
\item \st{L} \code{ibenarb/conway99-research} (dépôt actif, lu le 23/09) : même quotient $Q=2D_T+S+2L$ ; démontre $\tau\equiv0\pmod3$ par un comptage
élégant (les $231$ triangles du graphe ; les triangles fixés par $g$ sont exactement les $\tau$ orbites triangulaires) ; exclut les 4-cycles du 2-facteur $L$
et démontre le « lemme B » (une composante triangulaire $\{a,b,c\}$ de $L$ reçoit exactement une arête simple de chaque autre orbite).
Découpage par type de cycles de $L$ : 103 types « ouverts » de leur point de vue ; avec Ishida ($\tau=6$), leurs 2 types $\tau=27$ sont vides.
Coût mesuré de leur méthode : l'exclusion d'\emph{un} type, $(6,3^7)$, a demandé 656 sous-cas et 652\,Go de preuves LRAT.
\item \st{L} \code{Lundle22/conway-99-cyclic-lifts} : sous l'équation du quotient, l'équation complète équivaut à la congruence $W^2+W-12I\equiv0\pmod{\pi^3}$
($\pi=1-\omega$) sur la matrice d'Eisenstein, soit une équation linéaire et une quadratique sur $\F_3$ (lemme clé revérifié : trois entiers $\ge0$ de somme 6
avec $\pi^3\mid c_0+c_1\omega+c_2\omega^2$ valent tous 2). Utile pour l'étage des phases.
\item \st{L} Thakkar--Severini (arXiv:2608.11211) : modèle d'orbites générique ; ordre 3 sans point fixe indécis en 30\,min. SAT Modulo Symmetries :
non adapté (pas de multigraphe ni de partition imposée). « Smart cubing » : gains $\times2$--$3$.
\end{itemize}

\subsection{La relaxation linéaire LQ}
\begin{lemma}[bloc triangles $\times$ indépendantes] \st{D}\label{lem:LQ}
Pour toute orbite indépendante $Z$ de support $s_Z$ : $\sum_{Y\neq Z}C_{YZ}\,s_Y=t_Z:=6\cdot\mathbf1-3s_Z-Fs_Z$, $\sum_Y C_{YZ}=14-|s_Z|$, et exactement deux
$Y$ ont $C_{YZ}=2$ ; en particulier $s_Y(i)\,C_{YZ}\le t_Z(i)$.
\end{lemma}
\begin{proof}
Bloc $(i,Z)$ de $Q^2+Q=12I+6J$ : $3s_Z(i)+(Fs_Z)_i+\sum_Y s_Y(i)C_{YZ}=6$ ; somme de ligne de $Q$ ; lemme~\ref{lem:entries}(c).
\end{proof}
\paragraph{Résultats.} Sur les 33 problèmes quotient : \textbf{15 infaisables} pour LQ, dont les 8 déjà certifiés et \textbf{7 nouveaux}
(m04776\_s0000, s0002, s0003 ; m04905\_s0004, s0005 ; m05033\_s0008 ; m15426\_s0017). \textbf{13 des 15 le sont par une seule ligne} : une orbite $Z$ n'admet
aucun voisinage compatible ; certificat = recherche exhaustive sur une ligne (Python pur, \code{lq\_local.py}), recoupée par un modèle CP-SAT indépendant
sur les 891 lignes des 33 problèmes (0 désaccord). Les 2 autres (m16927\_s0018, s0023) demandent la symétrie globale et sont certifiés DRAT.
La variante la plus faible (sans lemme B, sans règle des 4-cycles, sans borne sur les doubles) suffit dans les 15 cas.

\paragraph{Contenu linéaire épuisé.} \st{D} Toute somme de $C^2$ contre un vecteur de $W=\operatorname{Im}[B^{\mathsf T}|\mathbf1]$ se réduit, via le
lemme~\ref{lem:LQ}, à une identité : LQ contient tout le contenu linéaire. Le reste est quadratique, dans $W^\perp$ ($C^2+C=12I$ y restreint).

\subsection{Classes clairsemées : énumération directe des supports}
Au lieu d'énumérer les $B$ avec leurs phases (explosif pour $F$ peu dense), on énumère les multi-ensembles de supports compatibles avec LQ, graphe par graphe
(\code{enum\_supports\_lq.py} : supports triés + $C$ + contraintes linéaires ; chaque solution est interdite puis on recommence ; fin sur « infaisable »
$\Rightarrow$ liste complète).
\begin{itemize}[nosep]
\item $F=K_{1,3}$ (masque 28, classe 42, auparavant 637 $B$ phasés sans fin d'énumération) : \textbf{1 seul multi-ensemble}, liste complète en 59\,s.
\item $F=3K_2$ (masque 16420, classe 27, auparavant 1\,817 $B$ phasés) : \textbf{4 multi-ensembles}, liste complète en 102\,s ; cohérent avec l'ancienne
liste partielle (les 2 autres supports qu'elle contenait échouent à LQ).
\item Masque 8290 (classe 25) : 2 multi-ensembles ; masque 9313 (classe 32) : 2 ; masque 13452 (classe 5) : 1 ; masque 29448 (classe 6) : 1.
Chaque fois, les multi-ensembles de l'ancienne liste partielle absents de la nouvelle échouent bien à LQ (\code{compare\_supports.py}).
\end{itemize}
Statut \st{S} : complétude affirmée par CP-SAT.

\paragraph{Test de phases sur un multi-ensemble de supports.} \code{support\_phase\_check.py} : existe-t-il un $B$ phasé de niveau 2.5 (résidu exact,
filtre de type, condition des triplets) d'une classe de $F$ dont les supports sont exactement ce multi-ensemble ? Si aucune classe ne le réalise,
le multi-ensemble est exclu. C'est valide avec les représentants de classes : tout graphe se réétiquette pour que sa classe soit un représentant, et son
multi-ensemble dans cet étiquetage est alors réalisé. \textbf{m08290\_sL1 et m09313\_sL1 sont exclus} ; certificat : vecteur de Farkas entier
(\code{farkas\_phase.py}), revérifié en arithmétique exacte \st{C}. Bilan des classes clairsemées : classe 42 $\to$ 1 problème quotient, 27 $\to$ 4,
25 $\to$ 1, 32 $\to$ 1, 5 $\to$ 1, 6 $\to$ 1.

\subsection{Accélération de l'étage Q : encodage compact et symétrie complète}
\paragraph{Encodage compact \st{D}+\st{C}.} L'ancien générateur encodait chaque contrainte « $\sum\ell=r$ » par \code{CardEnc.equals} (totalizer) ;
la moitié « $\ge r$ » y devient « au plus $n-r$ » sur les négations, soit un totalizer presque complet en $O(n^2)$ pour les $\approx100$ littéraux
d'une équation de Gram où $r\le6$. \code{lean\_encoding.py} utilise \emph{un} totalizer bidirectionnel dont les sorties sont tronquées à $r+1$
($o_j\Leftrightarrow\sum\ge j$), tout aussi arc-cohérent, en $O(nr^2)$ clauses ; variables et contraintes sont inchangées.
Validation (\code{validate\_lean.py}, PASS) : 5\,116 affectations exhaustives ($n\le8$, tous $r$) ; quotients réels BvLS (81 T et 81 I, lignes
libres : SAT et évalué ; entrée corrompue : UNSAT, pour les deux encodages) ; graphe tour $3\times3$ ; deux instances certifiées UNSAT.
\begin{center}\begin{tabular}{lrrrr}
instance & clauses (ancien) & clauses (compact) & temps ancien (natif) & temps compact\\\hline
m05033\_s0009 & 4\,727\,236 & 1\,117\,843 & 60\,s & 2{,}8\,s\\
m29481\_s0030 & 4\,721\,852 & 1\,122\,170 & 260\,s & 8{,}0\,s
\end{tabular}\end{center}

\paragraph{Symétrie complète \st{D}.} Jusqu'ici seules les orbites « jumelles » (supports identiques) étaient échangées. Une instance
$(F,\{s_Y\})$ admet aussi les permutations $\pi$ des 6 triangles avec $\pi(F)=F$ et $\pi(\{s_Y\})=\{s_Y\}$, prolongées aux orbites indépendantes :
71 non triviales pour m23221\_s0027, 23 pour m16420\_sL*, 11 pour m21041, m23431, m00028\_sL0. On impose $\mathrm{ch}(Q)\le_{\mathrm{lex}}\mathrm{ch}(Q^\sigma)$
pour chacune, où $\mathrm{ch}$ lit la matrice $33\times33$ complète ligne par ligne (ordre $0<1<2$).
\begin{lemma}\label{lem:symlex}
Soit $Q_0$ le minimum lexicographique de l'orbite d'une solution sous le groupe engendré par ces $\sigma$ et les transpositions de jumelles. Alors
$Q_0$ satisfait les contraintes $\sigma$ et les contraintes jumelles (ligne$_i\le_{\mathrm{lex}}$ligne$_j$ hors colonnes $i,j$, pour $i<j$ jumelles).
\end{lemma}
\begin{proof}
Contraintes $\sigma$ : minimalité. Jumelles : sinon, soit $c$ la première colonne ($\notin\{i,j\}$) où les lignes $i,j$ diffèrent, avec $Q_0[i][c]>Q_0[j][c]$.
Pour $k<c$, $k\notin\{i,j\}$ : $Q_0[k][i]=Q_0[i][k]=Q_0[j][k]=Q_0[k][j]$. Si $c<i$, la transposition $(ij)$ laisse les lignes $<c$ inchangées
et rend la ligne $c$ plus petite en colonne $i$. Si $c>i$, elle laisse les lignes $<i$ inchangées et la ligne $i$ diffère d'abord en $c$, plus petite
(les diagonales coïncident : même type). Contradiction avec la minimalité.
\end{proof}
Validation (\code{validate\_symlex.py}, PASS, avec les fonctions de production) : 60 petits modèles énumérés exhaustivement, groupe complet des données
fixées calculé par force brute ; 4\,044 orbites de solutions, \textbf{aucune perdue}, et le minimum lex de chaque orbite toujours présent ;
36 modèles ont des symétries $\sigma$ actives et des solutions (jusqu'à 13\,950 solutions, 3\,558 orbites). Témoin : la version relâchée de m23221\_s0027
(équations de Gram indépendantes retirées) reste SAT avec les 71 contraintes.

\paragraph{Mémoire.} Solveur alimenté clause par clause depuis le disque (\code{solve\_stream.py}) : 200--350\,Mo par instance au lieu de
$\approx1{,}4$\,Go, d'où 6 solveurs en parallèle au lieu de 2.

\subsection{Premier résultat de l'étage Q accéléré}
\textbf{m16420\_sL0 est UNSAT} \st{C} : CNF compacte avec symétrie complète (sha256 \code{a706468daf92ee33}), UNSAT en 2\,273\,s (CaDiCaL 1.9.5) puis
3\,515\,s (CaDiCaL 3.0.1 natif, autre version) ; preuve DRAT binaire de 129\,Mo vérifiée par drat-trim (2\,852\,s ; 587\,320 lemmes dans le noyau,
aucun lemme RAT). m16420\_sL3 lui est isomorphe. La validité repose en plus sur le lemme~\ref{lem:symlex} et sur la validation de l'encodage.

\subsection{Cube-and-conquer}
Les autres instances ne concluent pas en une heure (étage Q monolithique). Découpage : chaque orbite indépendante a exactement deux doubles
(lemme~\ref{lem:entries}(c)), donc pour une orbite $Z$ fixée, les 325 « cubes » $\{w_{ZA}=w_{ZB}=1\}$ couvrent toutes les solutions. Un cube
récalcitrant se redécoupe de la même façon : une orbite $V$ n'ayant qu'un double supposé (partenaire $P$) a son second double $W\notin\{V,P\}$,
d'où 25 sous-cubes qui couvrent le nœud (profondeur $\le2$). Résolution \emph{incrémentale} (\code{cube\_q\_inc.py}) : un seul solveur par
instance, chaque cube résolu sous hypothèses avec un budget de conflits (5\,000, puis 50\,000, puis découpage) ; les clauses apprises, conséquences
de la CNF de base, servent d'un cube à l'autre.
\paragraph{Validations.} \code{validate\_inc.py} (PASS) : sur le témoin relâché (m23221\_s0027 sans les équations de Gram indépendantes), 325 cubes,
verdicts identiques entre l'incrémental et un solveur neuf avec clauses unitaires (29 SAT, 296 UNSAT, modèles revérifiés clause par clause).
Le gain vient surtout de ce qu'on ne recharge plus 1,4\,M de clauses par cube : la plupart des cubes se résolvent en quelques millisecondes.
\code{validate\_split.py} (PASS) : sur les 60 derniers cubes du témoin avec un budget d'un conflit, 11 cubes tranchés par découpage (7 SAT trouvés
au fond de l'arbre, 1 UNSAT), 0 désaccord, 2\,819 nœuds.

\subsection{Doublons par isomorphisme}
Deux problèmes quotient $(F,\mu)$ et $(F,\pi(\mu))$ avec $\pi\in\operatorname{Aut}(F)$ ont la même réponse. Test systématique sur les 44 fichiers
d'instances : m04776\_s0000$\sim$s0002, m04905\_s0004$\sim$s0005, m16927\_s0018$\sim$s0023, m29481\_s0030$\sim$s0031 (morts des deux côtés),
m16927\_s0019$\sim$s0022, m16927\_s0020$\sim$s0021 (ouverts), m16420\_sL0$\sim$sL3, sL1$\sim$sL2. Les 33 instances d'origine ne font que
\textbf{27 classes d'isomorphisme}, et les 18 ouvertes \textbf{16}. Contrôle de cohérence : chaque paire a le même statut (LQ comme étage Q).

\subsection{Énumération des supports avec les phases}
\code{enum\_supports\_phased.py} : pour une classe $c$, un seul modèle CP-SAT contient le modèle de supports (LQ) et un $B$ phasé de niveau 2.5
de la classe, liés par « nombre de colonnes phasées de support $v$ = nombre d'orbites de support $v$ ». On bloque chaque multi-ensemble trouvé ;
« infaisable » à la fin donne la liste complète des multi-ensembles compatibles LQ \emph{et} réalisables par la classe. Validation
(\code{validate\_enum\_phased.py}, PASS) : sur les classes 0, 15, 30 (listes connues par l'énumération phasée complète) et 25, 32, 42
(clairsemées), ensembles identiques, listes complètes en 17--296\,s.
Variante négative \st{N} : un lex conditionnel des lignes de $C$ pour les supports égaux ralentit CP-SAT (masque 28 : rien en 600\,s au lieu de 135\,s).

\subsection{Pistes testées sans effet \st{N}}
\begin{itemize}[nosep]
\item LQ au niveau d'un graphe $F$ entier (supports libres) : aucun $F$ exclu (21 faisables, 5 indécis en 5\,min).
\item Version « phases » locale du lemme~\ref{lem:LQ} (cible $6\times3$ par orbite ; conventions validées sur 14\,580 comparaisons) : aucune élimination sur 15 classes.
\item Version « phases » globale (\code{phase\_global.py}, validée par contrôles plantés, jamais de faux négatif) : indécise en 60\,s sur les classes
symétriques 36--41.
\item Générateur ligne par ligne : 35\,000 à 2 millions de lignes possibles par orbite ; non viable sans élagage quadratique fort.
\item Balayage de la version globale sur 145 $B$ phasés à 15\,s chacun : 0 faisable, 0 exclu, tout indécis.
\item Relaxation LQB (LQ + inégalité de budget) : indécise en CP-SAT sur toutes les instances testées.
\item Condition spectrale par factorisation de $\chi_M(x^2+x)$ : déjà contenue dans le test spectral de l'audit ($CM=MK$ fixe tout le spectre de $C$).
\item Solutions invariantes (ansatz) pour m23221\_s0027 : aucune sous 8 types de symétries (6 UNSAT, 2 indécis à 900\,s). Rien n'indique que l'étage Q
y soit satisfiable.
\item Les lemmes « pas de $C_4$ dans $L$ » et « lemme B » sont \emph{impliqués} par l'étage Q (Gram $\ge8>6$, resp.\ Gram saturé puis sommes de lignes).
\end{itemize}

\subsection{Bilan au matin du 23 septembre (07\,h\,10) et suite}
\paragraph{Acquis de la nuit.}
\begin{itemize}[nosep]
\item \st{C} m16420\_sL0 (et son isomorphe sL3) : UNSAT, preuve DRAT vérifiée.
\item \st{S} m00028\_sL0 : UNSAT par cube-and-conquer (325/325 cubes, dont 27 fermés par découpage, 675 sous-cubes) ; c'était l'unique instance de la
classe 42 ($F=K_{1,3}$), qui est donc exclue à l'étage quotient. Certification DRAT cube par cube à faire.
\item Instances d'origine : 27 classes d'isomorphisme sur 33 ; ouvertes : 16 sur 18.
\item \st{S} Classe 4 (masque 30992) : liste complète avec phases, 96 multi-ensembles, 18 à isomorphisme près (m30992\_sP0--sP17), pas encore passées à l'étage Q.
\end{itemize}
\paragraph{Instances en cours (cubes UNSAT sur 325, aucun cube SAT).}
\begin{center}\begin{tabular}{lr|lr|lr}
m16420\_sL1 & 316 & m21041\_s0025 & 313 & m23221\_s0027 & 313\\
m23431\_s0028 & 309 & m06717\_s0014 & 307 & m20723\_s0024 & 306\\
m04907\_s0007 & 301 & m12295\_s0015 & 297 & m05033\_s0012 & 291\\
m29481\_s0032 & 288 & & & &
\end{tabular}\end{center}
Pas encore commencées : m04776\_s0001, m04905\_s0006, m15426\_s0016, m16927\_s0019, s0020, m21045\_s0026, m26850\_s0029, m08290\_sL0, m09313\_sL0,
m13452\_sL0, m29448\_sL0, m30992\_sP0--sP17.
\paragraph{Classes encore incomplètes.} 29 ($F=K_{1,5}$ : au moins 16 multi-ensembles avec phases, énumération en cours), 43 (2 trouvés, puis
indécis à 3\,600\,s), 44 et 45 (pas encore traitées avec les phases ; $F$ à une arête et $F$ vide, énumération sans phases indécise).
\paragraph{Coût restant.} Les passes à 5\,000 puis 50\,000 conflits règlent 85 à 95\,\% des cubes en quelques minutes ; chaque cube dur restant
coûte environ 30--45\,min de CPU par découpage. Il en reste de l'ordre de 300 à 400 sur les instances en cours, et les instances non commencées
s'y ajoutent : plusieurs centaines d'heures-CPU. C'est hors de portée de ce portable (mémoire : 15\,Go partagés avec Windows, trois alertes
mémoire cette nuit), mais adapté à la machine de calcul du projet (32\,Go) sur quelques jours, ou à un cluster.
\paragraph{Arrêt sur ce portable (23 septembre, 10\,h\,14).} À la demande de l'utilisateur, les solveurs de cubes ont été arrêtés proprement :
3\,381 cubes UNSAT au total, 50 cubes durs fermés par découpage, 1 indécis, aucun cube SAT. Par instance (UNSAT sur 325) : m00028\_sL0 325,
m16420\_sL1 324, m23221\_s0027 315, m21041\_s0025 313, m23431\_s0028 309, m06717\_s0014 307, m20723\_s0024 306, m12295\_s0015 302,
m04907\_s0007 301, m05033\_s0012 291, m29481\_s0032 288. L'énumération de la classe 29 a été arrêtée avec 31 multi-ensembles enregistrés (liste
incomplète).
\paragraph{Réénumération avec phases et comparaison.} \code{enum\_supports\_phased.py} a été relancé sur les classes de l'ancienne chaîne.
Toutes celles terminées ont exactement \textbf{un} multi-ensemble compatible, sauf la classe 0 (4) et la classe 4 (96). Comparaison avec
l'ancienne chaîne indépendante (énumération phasée de $B$ puis filtre LQ, \code{compare\_phased\_old.py}) sur 30 classes : 27 identiques ;
les 3 écarts sont attendus (classe 4 : l'ancienne liste ne contient que les 18 représentants par symétrie des 96 ; classes 25 et 32 : l'ancienne
liste contenait m08290\_sL1 et m09313\_sL1, exclus par les phases avec certificat de Farkas).

\paragraph{Pour reprendre.} Tout est reprenable : \code{cubes/LIST.txt} (file), \code{cubes/<instance>/results.txt} et \code{subcubes.txt}
(résultats), \code{sh cube\_runner\_inc.sh 5000} (un runner par cœur, environ 300\,Mo chacun), marqueur \code{FROM\_SPLIT} pour aller
directement au découpage. Certification : \code{cert\_sym.sh} (instance entière) ; pour les cubes, une CNF = base + clauses unitaires du cube.
